\documentclass[12pt]{article}
\usepackage[utf8]{inputenc}
\usepackage[french]{babel}
\usepackage{amsmath}
\usepackage[siunitx]{circuitikz}
\usepackage{pgfplots}
\usetikzlibrary{calc}

\begin{document}

\section*{Circuit RC et Charge d'un Condensateur}
\subsection*{Schéma du Circuit RC}
Voici un schéma d'un circuit RC en série :
\begin{center}
\begin{tikzpicture}[american]
    \node[ground] at (0,0) {};
    \draw (0,0) node[battery1, l=$\mathcal{V}_s$] {} -- ++(0,3);
    \draw (0,3) to[R, l=$R$] (0,0);
    \draw (0,3) to[C={i>\pgfkeysvalueof{/pgf/minimum width}}, l_=$C$] ++(3,0);
    \draw (3,3) -- (3,0) to[short] (0,0);
\end{tikzpicture}
\end{center}

\subsection*{Graphique de la Charge du Condensateur}
Voici la courbe de la tension aux bornes du condensateur en fonction du temps :

\begin{center}
\begin{tikzpicture}
\begin{axis}[
    xlabel={Temps $t$ (\si{\second})},
    ylabel={Tension $V_C(t)$ (\si{\volt})},
    xmin=0, xmax=0.01,
    ymin=0, ymax=15,
    samples=200,
    axis lines=middle,
    grid=major,
    width=10cm,
    height=6cm
]

% Paramètres du circuit
\def\VS{15} % Tension source en volts
\def\RC{0.005} % Produit RC en secondes

\addplot[domain=0:0.01, blue] { \VS * (1 - exp(-x/\RC)) };
\addlegendentry{ $\mathcal{V}_C(t) = \mathcal{V}_s (1 - e^{-t/RC})$ };

\end{axis}
\end{tikzpicture}
\end{center}

\end{document}